%% ===================================================================
%%  Doc Math DZ Book Series -- Volume 1: Linear Algebra
%%  Master file. Content lives in chapters/, figures in figures/.
%%  Build:  latexmk -pdf book.tex   (see .latexmkrc)
%% ===================================================================
\documentclass[print]{docmathdz}

%% ---- volume metadata (drives cover, title page and PDF properties) ----
\dmzset{
  title    = Linear Algebra,
  subtitle = Theory, Intuition and Doctoral Entrance Problems,
  subject  = Linear Algebra,
  tagline  = Theory \textperiodcentered\ Intuition \textperiodcentered\ Proofs
             \textperiodcentered\ Examples \textperiodcentered\ PhD Problems
             \textperiodcentered\ Solutions,
  author   = Abdelouahab Mostafa,
  volume   = 1,
  accent   = dmzblue,
  edition  = First Edition,
  year     = 2026,
}

\addbibresource{references.bib}

\begin{document}

%% ---------------- front matter ----------------
\dmzfrontmatter{%
  Volumes of the series:\par\medskip
  \begin{tabular}{@{}r@{\quad}l@{}}
    1 & Linear Algebra\\
    2 & Real Analysis\\
    3 & Topology\\
    4 & Measure Theory\\
    5 & Functional Analysis\\
    6 & Probability Theory\\
    7 & Group Theory\\
    8 & Ring Theory\\
    9 & Field Theory\\
   10 & Complex Analysis\\
   11 & Differential Equations\\
   12 & Dynamical Systems\\
   13 & Numerical Analysis\\
   14 & Optimization\\
   15 & Graph Theory\\
  \end{tabular}\par\medskip
  Each volume is devoted to a single subject and can be studied independently.
}

\chapter*{Preface}
\addcontentsline{toc}{chapter}{Preface}
\markboth{Preface}{Preface}

This book was written for one precise situation: you are a Master's student in
Algeria, you intend to sit a doctoral entrance examination, and you want to study
linear algebra seriously, alone, with limited time.

Three decisions follow from that situation. First, the book covers \emph{one}
subject only; nothing here belongs to analysis or probability. Second, every
notion is introduced twice: once informally, so that you see what it means, and
once formally, so that you can prove things with it. Third, the theory is
connected to real examinations: the last two chapters of the book are entirely
devoted to doctoral entrance problems and to their complete solutions.

Nothing is left as an exercise without a solution. If a step is routine it is
still written down, because a step that is routine for the author is rarely
routine for a reader meeting it for the first time.

\bigskip\noindent\emph{Abdelouahab Mostafa}\\
\emph{Doc Math DZ}

\chapter*{Learning Objectives}
\addcontentsline{toc}{chapter}{Learning Objectives}
\markboth{Learning Objectives}{Learning Objectives}

After working through this volume you should be able to:

\begin{enumerate}
  \item decide whether a given set with two operations is a vector space, and
        recognise the standard examples and counter-examples;
  \item manipulate families of vectors fluently: freeness, generation, bases,
        dimension, and the exchange arguments behind them;
  \item compute the rank of a linear map or of a matrix by three independent
        methods and explain why they agree;
  \item use the rank--nullity theorem as a counting tool rather than as a formula;
  \item determine whether an endomorphism is diagonalisable, and diagonalise it
        explicitly when it is;
  \item exploit annihilating polynomials, in particular the criterion by split
        polynomials with simple roots;
  \item state and apply the spectral theorem for real symmetric matrices;
  \item solve a complete doctoral entrance problem in linear algebra within the
        time of an examination.
\end{enumerate}

\chapter*{Roadmap of the Subject}
\addcontentsline{toc}{chapter}{Roadmap of the Subject}
\markboth{Roadmap}{Roadmap}

Linear algebra has a very short logical skeleton, and it pays to keep it in view.

\begin{center}
\begin{tikzpicture}[
    node distance=7mm and 10mm,
    every node/.style={font=\sffamily\small},
    box/.style={draw=dmzaccent,fill=dmzaccent!5,rounded corners=1.5pt,
                inner sep=5pt,text width=32mm,align=center},
    arr/.style={-{Stealth[length=2mm]},draw=dmzsoft}]
  \node[box] (vs) {Vector spaces};
  \node[box,below=of vs] (fam) {Families: free, generating, bases};
  \node[box,below=of fam] (dim) {Dimension};
  \node[box,right=of dim] (maps) {Linear maps};
  \node[box,right=of maps] (rank) {Rank and rank--nullity};
  \node[box,below=of rank] (mat) {Matrices, change of basis};
  \node[box,below=of mat] (red) {Reduction: eigenvalues, diagonalisation};
  \node[box,left=of red] (spec) {Spectral theorem};
  \draw[arr] (vs) -- (fam);
  \draw[arr] (fam) -- (dim);
  \draw[arr] (dim) -- (maps);
  \draw[arr] (maps) -- (rank);
  \draw[arr] (rank) -- (mat);
  \draw[arr] (mat) -- (red);
  \draw[arr] (red) -- (spec);
\end{tikzpicture}
\end{center}

Everything else in the subject is a refinement of one of these eight boxes.

\chapter*{Prerequisites}
\addcontentsline{toc}{chapter}{Prerequisites}
\markboth{Prerequisites}{Prerequisites}

The reader is assumed to know: elementary set theory and maps (injective,
surjective, bijective); the notion of a field, with $\Q$, $\R$, $\C$ and
$\Z/p\Z$ as examples; polynomial arithmetic in $\K[X]$, including division and
factorisation; and Gaussian elimination on a system of linear equations.

Nothing from analysis is used, except in a few remarks that are explicitly
marked as such.

%% ---------------- main matter ----------------
\mainmatter

\chapter{Introduction and Orientation}
\label{ch:introduction}

\section{Motivation}

Consider a system of linear equations
\begin{equation}\label{eq:system}
  \left\{\;\begin{aligned}
    a_{11}x_1+\dots+a_{1n}x_n &= b_1,\\
                              &\;\;\vdots\\
    a_{m1}x_1+\dots+a_{mn}x_n &= b_m.
  \end{aligned}\right.
\end{equation}
Every student can solve such a system by elimination. But elimination answers
only the question \emph{what are the solutions?}. It does not answer the
questions an examiner actually asks: for how many right-hand sides $b$ does a
solution exist? When is it unique? If the data $a_{ij}$ are perturbed slightly,
does the answer change? What is common to \eqref{eq:system}, to a linear
differential equation, and to a recurrence relation?

Linear algebra is the theory that replaces the computation by a structural
answer. The set of solutions of \eqref{eq:system} with $b=0$ is a vector space;
its dimension is $n-\rank A$; the system with $b\neq0$ is solvable exactly when
$b$ lies in the image of the associated map. Three notions -- space, map, rank --
absorb the whole question.

\begin{intuition}
Linear algebra studies the transformations of space that send straight lines to
straight lines and fix the origin. Such a transformation may stretch, rotate,
shear or flatten space, but it never bends it. Because the deformation is
uniform, the whole transformation is determined by what it does to finitely many
vectors: that is why the subject is computable, and also why it is everywhere.
\end{intuition}

\section{Historical background}

Determinants appear before matrices: Leibniz uses them in 1693, Cramer publishes
his rule in 1750. The matrix as an object in its own right is due to Cayley
(1858), who also states the theorem that bears his name with Hamilton. The
axiomatic notion of a vector space, freed from coordinates, is much later: it is
made explicit by Peano (1888) and becomes standard only with the treatises of
the 1930s, notably van der Waerden and Banach for the infinite-dimensional case.
The modern teaching order -- spaces first, matrices as a representation -- thus
reverses the historical order. Keep this in mind: coordinates were the beginning
of the subject, not its content.

\section{How to use this book}

The volume follows the fixed structure of the series.

\begin{itemize}
  \item \Cref{ch:definitions} contains the core definitions with their examples
        and, just as importantly, their non-examples.
  \item \Cref{ch:theorems} contains the fundamental theorems with complete proofs.
  \item \Cref{ch:examples} works out detailed computations, with figures.
  \item \Cref{ch:summary} gathers the common mistakes, a summary, and the
        practice exercises.
  \item \Cref{ch:problems} contains doctoral entrance problems, and
        \cref{ch:solutions} their fully worked solutions, together with the
        solutions of the exercises.
\end{itemize}

A reasonable working method: read a section, close the book, reconstruct the
statements from memory, and only then look at the proofs. Solve the exercises of
\cref{ch:summary} before reading any solution; a solution read is worth roughly
a tenth of a solution found.

\begin{note}[Notation]
Throughout, $\K$ denotes an arbitrary field, $E$ and $F$ are $\K$-vector spaces,
and $u,v$ denote linear maps. $\Mat{n}{\K}$ is the algebra of $n\times n$
matrices over $\K$, $\id$ the identity map, $\transpose{A}$ the transpose of $A$,
and $\spec(u)$ the set of eigenvalues of $u$.
\end{note}

\section{Further reading}

For a coordinate-free treatment close in spirit to this volume, see
\textcite{axler2015}. For the matrix-analytic point of view, \textcite{horn2013}
is the standard reference. Readers preparing Algerian doctoral examinations will
recognise the style of exercises in \textcite{gourdon2009} and
\textcite{grifone2011}.

\chapter{Core Definitions}
\label{ch:definitions}

\section{Motivation}

We want a framework in which \emph{addition} and \emph{scaling} make sense and
nothing else is assumed. The gain is uniformity: one theorem proved in that
framework applies simultaneously to tuples of numbers, to polynomials, to
matrices and to functions.

\section{Vector spaces}

\begin{definition}{Vector space}{vector-space}
Let $\K$ be a field. A \term{vector space} over $\K$ is a set $E$ equipped with
an addition $E\times E\to E$ and a scalar multiplication $\K\times E\to E$ such
that $(E,+)$ is an abelian group and, for all $\lambda,\mu\in\K$ and
$x,y\in E$,
\[
  \lambda(x+y)=\lambda x+\lambda y,\qquad
  (\lambda+\mu)x=\lambda x+\mu x,\qquad
  (\lambda\mu)x=\lambda(\mu x),\qquad
  1\cdot x=x.
\]
The elements of $E$ are called vectors, those of $\K$ scalars.
\end{definition}

\begin{example}{The four standard examples}{standard-spaces}
\begin{enumerate}
  \item $\K^{n}$ with componentwise operations.
  \item $\K[X]$, and its subspace $\K_{n}[X]$ of polynomials of degree at most $n$.
  \item $\Mat{m,n}{\K}$, the $m\times n$ matrices.
  \item $\mathcal{F}(X,\K)$, all maps from a set $X$ to $\K$.
\end{enumerate}
Every example met in this book is one of these four, or a subspace of one of them.
\end{example}

\begin{example}{Three non-examples}{non-examples}
\begin{enumerate}
  \item $\N^{2}$ is not a vector space over $\R$: it is not stable under
        multiplication by $-1$.
  \item $\set{(x,y)\in\R^{2}\setst xy=0}$ is stable under scaling but not under
        addition: $(1,0)+(0,1)=(1,1)$.
  \item The set of polynomials of degree \emph{exactly} $n$ is not a subspace:
        it does not contain $0$, and $X^{n}+(-X^{n}+1)$ has degree $0$.
\end{enumerate}
\end{example}

\begin{definition}{Subspace}{subspace}
A subset $F\subseteq E$ is a \term{linear subspace} if $0\in F$ and
$\lambda x+\mu y\in F$ for all $x,y\in F$ and $\lambda,\mu\in\K$. A subspace is
itself a vector space for the restricted operations.
\end{definition}

\begin{proposition}{Intersections and sums}{sum-intersection}
If $F$ and $G$ are subspaces of $E$, then $F\cap G$ and
$F+G=\set{x+y\setst x\in F,\ y\in G}$ are subspaces, and $F+G$ is the smallest
subspace containing $F\cup G$. In general $F\cup G$ is \emph{not} a subspace.
\end{proposition}

\begin{proof}
Stability of $F\cap G$ and of $F+G$ is immediate. Any subspace $H$ containing
$F\cup G$ contains every sum $x+y$ with $x\in F$, $y\in G$, hence contains
$F+G$; and $F+G$ contains $F\cup G$. The last assertion is \cref{exo:union}.
\end{proof}

\begin{definition}{Direct sum}{direct-sum}
The sum $F+G$ is \term{direct}, written $F\oplus G$, if $F\cap G=\set{0}$.
Equivalently, every $x\in F+G$ has a unique decomposition $x=y+z$ with $y\in F$
and $z\in G$.
\end{definition}

\begin{intuition}
A direct sum is a coordinate system: to say $E=F\oplus G$ is to say that each
vector of $E$ has exactly one address, one component in $F$ and one in $G$.
Uniqueness is what makes projections well defined.
\end{intuition}

\section{Families of vectors}

\begin{definition}{Linear combination, span}{span}
Let $(x_i)_{i\in I}$ be a family of vectors of $E$. Its \term{span}, written
$\Span(x_i)_{i\in I}$, is the set of finite linear combinations
$\sum_{i\in J}\lambda_i x_i$ with $J\subseteq I$ finite and $\lambda_i\in\K$. It
is the smallest subspace containing the family.
\end{definition}

\begin{definition}{Free family, generating family, basis}{basis}
A family $(x_i)_{i\in I}$ is
\begin{itemize}
  \item \term{free} (or linearly independent) if every finite relation
        $\sum_{i\in J}\lambda_i x_i=0$ forces $\lambda_i=0$ for all $i\in J$;
  \item \term{generating} if $\Span(x_i)_{i\in I}=E$;
  \item a \term{basis} if it is both free and generating.
\end{itemize}
\end{definition}

\begin{proposition}{Characterisation of a basis}{basis-unique}
A family $\mathcal{B}=(e_i)_{i\in I}$ is a basis of $E$ if and only if every
vector of $E$ is a linear combination of $\mathcal{B}$ in exactly one way.
\end{proposition}

\begin{proof}
If $\mathcal{B}$ is a basis, existence follows from generation, and if
$\sum\lambda_i e_i=\sum\mu_i e_i$ then $\sum(\lambda_i-\mu_i)e_i=0$, so
$\lambda_i=\mu_i$ by freeness. Conversely, existence of a decomposition gives
generation, and uniqueness applied to $0=\sum 0\cdot e_i$ gives freeness.
\end{proof}

\begin{definition}{Dimension}{dimension}
If $E$ admits a finite basis, all its bases have the same cardinality
(\cref{thm:dimension-well-defined}); this common number is the \term{dimension}
$\dim E$. Otherwise $E$ is said to be infinite-dimensional.
\end{definition}

\section{Linear maps}

\begin{definition}{Linear map}{linear-map}
A map $u\colon E\To F$ is \term{linear} if $u(\lambda x+\mu y)=\lambda u(x)+\mu u(y)$
for all $x,y\in E$ and $\lambda,\mu\in\K$. The set of such maps is written
$\Hom(E,F)$, and $\End(E)=\Hom(E,E)$. An \term{endomorphism} is an element of
$\End(E)$, an \term{isomorphism} is a bijective linear map.
\end{definition}

\begin{definition}{Kernel, image, rank}{kernel-image}
For $u\in\Hom(E,F)$ one sets
\[
  \Ker u=\set{x\in E\setst u(x)=0},\qquad
  \im u=\set{u(x)\setst x\in E},\qquad
  \rank u=\dim\im u .
\]
$\Ker u$ is a subspace of $E$ and $\im u$ a subspace of $F$.
\end{definition}

\begin{proposition}{Injectivity and the kernel}{injective-kernel}
A linear map $u$ is injective if and only if $\Ker u=\set{0}$.
\end{proposition}

\begin{proof}
If $u$ is injective and $u(x)=0=u(0)$ then $x=0$. Conversely, if
$\Ker u=\set{0}$ and $u(x)=u(y)$, then $u(x-y)=0$, so $x-y=0$.
\end{proof}

\begin{definition}{Projection}{projection}
An endomorphism $p\in\End(E)$ with $p^{2}=p$ is called a \term{projection}.
\end{definition}

\Cref{fig:projection} shows the geometric content of the identity $p^2=p$:
projecting a second time changes nothing.

\begin{figure}[htbp]
  \centering
  % Orthogonal projection onto a line: the geometric picture behind p^2 = p.
\begin{tikzpicture}[scale=1.05,>=Stealth,font=\small]
  \draw[dmzrule,line width=0.4pt] (-0.6,0) -- (4.4,0);
  \draw[dmzrule,line width=0.4pt] (0,-0.6) -- (0,3.1);
  % the line D = im p
  \draw[dmzaccent,line width=1pt] (-0.4,-0.2) -- (4.2,2.1) node[right,text=dmzaccent] {$D=\im p$};
  % vector x and its projection
  \draw[->,dmzink,line width=0.9pt] (0,0) -- (1.6,2.6) node[above left] {$x$};
  \draw[->,dmztheorem,line width=0.9pt] (0,0) -- (2.65,1.33) node[below right,text=dmztheorem] {$p(x)$};
  \draw[dmzsoft,dashed] (1.6,2.6) -- (2.65,1.33);
  \draw[dmzsoft,line width=0.4pt] (2.65,1.33) ++(-0.18,0.09) -- ++(0.09,0.18) -- ++(0.18,-0.09);
  % the kernel direction
  \draw[->,dmzgreen,line width=0.8pt] (0,0) -- (-1.05,1.27) node[above,text=dmzgreen] {$\Ker p$};
  \node[dmzsoft,anchor=north west,font=\footnotesize,text=dmzsoft] at (0.15,-0.15) {$0$};
\end{tikzpicture}

  \caption{A projection $p$ decomposes the plane as $\Ker p\oplus\im p$.}
  \label{fig:projection}
\end{figure}

\section{Eigenvalues}

\begin{definition}{Eigenvalue, eigenvector, spectrum}{eigenvalue}
Let $u\in\End(E)$. A scalar $\lambda\in\K$ is an \term{eigenvalue} of $u$ if
$\Ker(u-\lambda\id)\neq\set{0}$; a nonzero element of that kernel is an
\term{eigenvector}, and $E_\lambda=\Ker(u-\lambda\id)$ is the associated
\term{eigenspace}. The set of eigenvalues is the \term{spectrum} $\spec(u)$.
\end{definition}

\begin{definition}{Diagonalisable endomorphism}{diagonalisable}
An endomorphism $u$ of a finite-dimensional space $E$ is \term{diagonalisable}
if $E$ admits a basis of eigenvectors of $u$; equivalently, if
$E=\bigoplus_{\lambda\in\spec(u)}E_\lambda$.
\end{definition}

\begin{definition}{Annihilating polynomial}{annihilating}
A polynomial $P\in\K[X]$ \term{annihilates} $u$ if $P(u)=0$. The monic
generator of the ideal of annihilating polynomials is the \term{minimal
polynomial} $\pi_u$.
\end{definition}

\begin{warning}
A definition is not a criterion. \Cref{def:diagonalisable} says what
diagonalisable means; it gives no method. The usable criteria are
\cref{thm:diagonalisable-criterion} and its corollary on split annihilating
polynomials with simple roots.
\end{warning}

\chapter{Fundamental Theorems and Proofs}
\label{ch:theorems}

\section{Bases and dimension}

\begin{theorem}{Exchange lemma (Steinitz)}{steinitz}
Let $E$ be a vector space, $L=(x_1,\dots,x_p)$ a free family and
$G=(y_1,\dots,y_q)$ a generating family of $E$. Then $p\leq q$, and one may
complete $L$ into a basis of $E$ using vectors of $G$.
\end{theorem}

\begin{proof}
We prove by induction on $k\in\set{0,\dots,p}$ that, after renumbering $G$, the
family $(x_1,\dots,x_k,y_{k+1},\dots,y_q)$ is generating. For $k=0$ this is the
hypothesis on $G$. Assume it for some $k<p$. Since the family
$(x_1,\dots,x_k,y_{k+1},\dots,y_q)$ generates $E$, we may write
\[
  x_{k+1}=\sum_{i\leq k}\alpha_i x_i+\sum_{j>k}\beta_j y_j .
\]
If all $\beta_j$ vanished, $x_{k+1}$ would be a combination of $x_1,\dots,x_k$,
contradicting the freeness of $L$; in particular $k<q$, which already gives
$p\leq q$. Choose $j_0>k$ with $\beta_{j_0}\neq0$ and renumber so that $j_0=k+1$.
Then
\[
  y_{k+1}=\beta_{k+1}^{-1}\Bigl(x_{k+1}-\sum_{i\leq k}\alpha_i x_i-\sum_{j>k+1}\beta_j y_j\Bigr),
\]
so $y_{k+1}$ belongs to the span of $(x_1,\dots,x_{k+1},y_{k+2},\dots,y_q)$,
which is therefore generating. This completes the induction. For $k=p$ we obtain
a generating family $(x_1,\dots,x_p,y_{p+1},\dots,y_q)$ containing $L$; deleting
one by one the vectors $y_j$ that are combinations of the previous ones yields a
basis containing $L$.
\end{proof}

\begin{theorem}{Dimension is well defined}{dimension-well-defined}
In a vector space admitting a finite generating family, all bases are finite and
have the same number of elements.
\end{theorem}

\begin{proof}
Let $\mathcal{B}$ and $\mathcal{B}'$ be two bases, with $\mathcal{B}'$ finite of
cardinality $q$ (such a basis exists by extracting one from a finite generating
family, as at the end of the previous proof). Every finite subfamily of
$\mathcal{B}$ is free, hence has at most $q$ elements by \cref{thm:steinitz};
therefore $\mathcal{B}$ is finite, say of cardinality $p$, and $p\leq q$.
Exchanging the roles of $\mathcal{B}$ and $\mathcal{B}'$ gives $q\leq p$.
\end{proof}

\begin{corollary}{}{dim-subspace}
If $\dim E=n<\infty$ and $F$ is a subspace of $E$, then $\dim F\leq n$, with
equality if and only if $F=E$. Moreover every free family of $n$ vectors and
every generating family of $n$ vectors is a basis.
\end{corollary}

\begin{proof}
A free family of $F$ is free in $E$, so has at most $n$ elements by
\cref{thm:steinitz}; a maximal one is a basis of $F$, whence $\dim F\leq n$. If
$\dim F=n$, a basis of $F$ is a free family of $n$ vectors of $E$; completing it
into a basis of $E$ (\cref{thm:steinitz}) adds nothing, so it is a basis of $E$
and $F=E$. The last statement follows by the same completion and extraction
arguments.
\end{proof}

\section{Rank}

\begin{theorem}{Rank--nullity}{rank-nullity}
Let $u\in\Hom(E,F)$ with $\dim E=n<\infty$. Then
\[
  \dim E=\rank u+\dim\Ker u .
\]
\end{theorem}

\begin{proof}
Let $(e_1,\dots,e_k)$ be a basis of $\Ker u$ and complete it into a basis
$(e_1,\dots,e_n)$ of $E$ by \cref{thm:steinitz}. We claim that
$(u(e_{k+1}),\dots,u(e_n))$ is a basis of $\im u$.

It is generating: for $x=\sum_{i=1}^{n}\lambda_i e_i$ we get
$u(x)=\sum_{i>k}\lambda_i u(e_i)$, because $u(e_i)=0$ for $i\leq k$.

It is free: if $\sum_{i>k}\lambda_i u(e_i)=0$, then
$z=\sum_{i>k}\lambda_i e_i\in\Ker u$, so $z=\sum_{i\leq k}\mu_i e_i$ for some
scalars $\mu_i$. The family $(e_1,\dots,e_n)$ being free, all $\lambda_i$ (and
all $\mu_i$) vanish.

Hence $\rank u=n-k=\dim E-\dim\Ker u$.
\end{proof}

\begin{corollary}{}{iso-criterion}
If $\dim E=\dim F<\infty$ and $u\in\Hom(E,F)$, the following are equivalent:
$u$ is injective; $u$ is surjective; $u$ is an isomorphism.
\end{corollary}

\begin{proof}
By \cref{thm:rank-nullity}, $\Ker u=\set{0}$ if and only if $\rank u=\dim E=\dim F$,
which by \cref{cor:dim-subspace} holds if and only if $\im u=F$.
\end{proof}

\begin{proposition}{Rank inequalities}{rank-inequalities}
For $u,v$ composable linear maps between finite-dimensional spaces,
\[
  \rank(u\circ v)\leq\min(\rank u,\rank v),\qquad
  \rank(u+v)\leq\rank u+\rank v .
\]
\end{proposition}

\begin{proof}
Since $\im(u\circ v)=u(\im v)\subseteq\im u$, we get $\rank(u\circ v)\leq\rank u$;
and $u$ restricted to $\im v$ has rank at most $\dim\im v=\rank v$. For the sum,
$\im(u+v)\subseteq\im u+\im v$ and $\dim(\im u+\im v)\leq\rank u+\rank v$.
\end{proof}

\section{Reduction of endomorphisms}

\begin{theorem}{Eigenspaces are independent}{eigen-independent}
Eigenvectors associated with pairwise distinct eigenvalues form a free family;
equivalently, the sum $\sum_{\lambda\in\spec(u)}E_\lambda$ is direct.
\end{theorem}

\begin{proof}
Suppose, for a contradiction, that some nontrivial relation
$\sum_{i=1}^{r}x_i=0$ holds with $x_i\in E_{\lambda_i}\setminus\set{0}$ and the
$\lambda_i$ pairwise distinct, and take $r$ minimal. Applying $u-\lambda_r\id$
kills $x_r$ and gives $\sum_{i=1}^{r-1}(\lambda_i-\lambda_r)x_i=0$, a relation
with $r-1$ nonzero terms since $\lambda_i\neq\lambda_r$. This contradicts
minimality (and for $r=1$ it contradicts $x_1\neq0$).
\end{proof}

\begin{theorem}{Diagonalisability criterion}{diagonalisable-criterion}
Let $\dim E=n<\infty$ and $u\in\End(E)$. Then $u$ is diagonalisable if and only
if $\sum_{\lambda\in\spec(u)}\dim E_\lambda=n$.
\end{theorem}

\begin{proof}
If $u$ is diagonalisable, a basis of eigenvectors splits into bases of the
eigenspaces, so the dimensions add up to $n$. Conversely, by
\cref{thm:eigen-independent} the sum of the eigenspaces is direct, hence of
dimension $\sum_\lambda\dim E_\lambda=n$, so it equals $E$ by
\cref{cor:dim-subspace}; concatenating bases of the eigenspaces gives a basis of
eigenvectors.
\end{proof}

\begin{theorem}{Split annihilating polynomial with simple roots}{split-simple}
Let $u\in\End(E)$ with $\dim E<\infty$. If there is a polynomial
$P=\prod_{i=1}^{r}(X-\lambda_i)$ with pairwise distinct $\lambda_i\in\K$ such
that $P(u)=0$, then $u$ is diagonalisable and
$E=\bigoplus_{i=1}^{r}\Ker(u-\lambda_i\id)$.
\end{theorem}

\begin{proof}
Set $P_i=\prod_{j\neq i}(X-\lambda_j)$. The polynomials $P_1,\dots,P_r$ have no
common root, hence $\gcd(P_1,\dots,P_r)=1$ and B\'ezout gives $Q_i\in\K[X]$ with
$\sum_i Q_iP_i=1$. Evaluating at $u$,
\[
  \id=\sum_{i=1}^{r}Q_i(u)P_i(u).
\]
For $x\in E$ put $x_i=Q_i(u)P_i(u)(x)$, so that $x=\sum_i x_i$. Since
$(X-\lambda_i)P_i=P$ and $P(u)=0$, we get $(u-\lambda_i\id)(x_i)=0$, that is
$x_i\in\Ker(u-\lambda_i\id)$. Hence $E=\sum_i\Ker(u-\lambda_i\id)$, and the sum
is direct by \cref{thm:eigen-independent}. Any concatenation of bases of these
kernels is a basis of eigenvectors.
\end{proof}

\begin{note}
\Cref{thm:split-simple} is the single most profitable theorem of this chapter in
an examination: an identity such as $u^{2}=u$, $u^{2}=\id$ or $u^{3}=u$
immediately produces diagonalisability, with no computation of a characteristic
polynomial.
\end{note}

\begin{theorem}{Spectral theorem, real symmetric case}{spectral}
Let $A\in\Mat{n}{\R}$ be symmetric. Then $A$ is diagonalisable in an orthonormal
basis: there exist an orthogonal matrix $Q$ and a real diagonal matrix $D$ with
$A=QD\transpose{Q}$. In particular $\spec(A)\subseteq\R$ and eigenspaces
associated with distinct eigenvalues are orthogonal.
\end{theorem}

\begin{proof}
First, the eigenvalues are real: if $Az=\lambda z$ with $z\in\C^{n}$ nonzero,
then $\adjoint{z}Az=\lambda\adjoint{z}z$ and, taking conjugate transposes and
using $\transpose{A}=A$ with real entries, $\adjoint{z}Az=\bar\lambda\adjoint{z}z$;
since $\adjoint{z}z>0$, $\lambda=\bar\lambda$.

Now argue by induction on $n$, the case $n=1$ being trivial. The characteristic
polynomial of $A$ has a root $\lambda_1\in\R$ by the previous paragraph and the
fundamental theorem of algebra, with a real eigenvector $q_1$, normalised so
that $\norm{q_1}=1$. Let $H=q_1^{\perp}$. For $x\in H$,
\[
  \inner{Ax}{q_1}=\inner{x}{Aq_1}=\lambda_1\inner{x}{q_1}=0,
\]
so $A(H)\subseteq H$. The restriction of $A$ to $H$ is symmetric for the induced
inner product, and $\dim H=n-1$, so by induction $H$ has an orthonormal basis
$(q_2,\dots,q_n)$ of eigenvectors. Then $(q_1,\dots,q_n)$ is an orthonormal
basis of $\R^{n}$ made of eigenvectors of $A$, and $Q=(q_1|\cdots|q_n)$ is
orthogonal with $\transpose{Q}AQ=D$ diagonal.
\end{proof}

\begin{corollary}{}{spd-eigen}
A symmetric matrix $A\in\Mat{n}{\R}$ is positive definite if and only if all its
eigenvalues are positive; in that case $A$ is invertible and $A^{-1}$ is
symmetric positive definite with eigenvalues the inverses of those of $A$.
\end{corollary}

\begin{proof}
Write $A=QD\transpose{Q}$ as in \cref{thm:spectral} and put $y=\transpose{Q}x$;
then $\transpose{x}Ax=\sum_i\lambda_i y_i^{2}$, which is positive for all
$x\neq0$ exactly when every $\lambda_i>0$. Then $D$ is invertible and
$A^{-1}=QD^{-1}\transpose{Q}$.
\end{proof}

\chapter{Illustrative Examples}
\label{ch:examples}

The theorems of \cref{ch:theorems} become usable only after they have been seen
at work. Every computation below is written in full, in the order in which it
should be produced on an examination sheet.

\section{Rank computed in three ways}

\begin{example}{One matrix, three methods}{rank-three-ways}
Let
\[
  A=\begin{pmatrix} 1 & 2 & 3\\ 2 & 4 & 6\\ 1 & 1 & 1 \end{pmatrix}
  \in\Mat{3}{\R}.
\]

\textbf{By elimination.} Subtracting twice the first row from the second gives a
zero row; subtracting the first row from the third gives $(0,-1,-2)$. Hence
\[
  A\sim\begin{pmatrix} 1 & 2 & 3\\ 0 & -1 & -2\\ 0 & 0 & 0\end{pmatrix},
  \qquad \rank A=2 .
\]

\textbf{By the image.} The columns are $c_1=(1,2,1)$, $c_2=(2,4,1)$,
$c_3=(3,6,1)$. One checks $c_3=2c_2-c_1$, while $c_1,c_2$ are not proportional,
so $\im A=\Span(c_1,c_2)$ has dimension $2$.

\textbf{By the kernel.} Solving $Ax=0$: the second equation is twice the first,
and the system reduces to $x_1+2x_2+3x_3=0$, $x_2+2x_3=0$, whence
$x_2=-2x_3$ and $x_1=x_3$. Thus $\Ker A=\Span\bigl((1,-2,1)\bigr)$ has
dimension $1$, and $\rank A=3-1=2$ by \cref{thm:rank-nullity}.

The three methods must agree, and comparing them is the cheapest way to detect
an arithmetic slip.
\end{example}

\section{A projection, seen algebraically and geometrically}

\begin{example}{Decomposition induced by a projection}{projection-example}
Let $p\colon\R^{2}\To\R^{2}$ be given in the canonical basis by
\[
  P=\frac{1}{5}\begin{pmatrix} 4 & 2\\ 2 & 1\end{pmatrix}.
\]
Then $P^{2}=\frac{1}{25}\begin{pmatrix}20&10\\10&5\end{pmatrix}=P$, so $p$ is a
projection. Its image is $\Span\bigl((2,1)\bigr)$ and its kernel is
$\Span\bigl((1,-2)\bigr)$, since $P(1,-2)^{\mathsf{T}}=0$. These two lines are
orthogonal, so $p$ is the orthogonal projection onto the line $y=x/2$, exactly
the situation of \cref{fig:projection}.

By \cref{thm:split-simple} applied to $X^{2}-X=X(X-1)$, the map $p$ is
diagonalisable with $\spec(p)=\set{0,1}$; in the basis $\bigl((2,1),(1,-2)\bigr)$
its matrix is $\diag(1,0)$. Note also $\tr P=1=\rank P$, an instance of
\cref{prob:idempotent}.
\end{example}

\section{Diagonalisation without a characteristic polynomial}

\begin{example}{An involution}{involution}
Let $u\in\End(\R^{3})$ with matrix
\[
  M=\begin{pmatrix} 0 & 1 & 0\\ 1 & 0 & 0\\ 0 & 0 & -1\end{pmatrix}.
\]
A direct computation gives $M^{2}=I_3$, so $X^{2}-1=(X-1)(X+1)$ annihilates $u$.
Its roots are simple, hence $u$ is diagonalisable by \cref{thm:split-simple} and
\[
  \R^{3}=\Ker(u-\id)\oplus\Ker(u+\id)
        =\Span\bigl((1,1,0)\bigr)\oplus\Span\bigl((1,-1,0),(0,0,1)\bigr).
\]
No determinant was computed. This is the standard time-saving move in an
examination.
\end{example}

\section{When diagonalisation fails}

\begin{example}{A non-diagonalisable matrix}{jordan-block}
Let $N=\begin{pmatrix}0&1\\0&0\end{pmatrix}$. Its only eigenvalue is $0$, and
$\Ker N=\Span\bigl((1,0)\bigr)$ has dimension $1<2$, so $N$ is not
diagonalisable by \cref{thm:diagonalisable-criterion}. Observe that $X^{2}$
annihilates $N$: an annihilating polynomial that is split but has a
\emph{multiple} root gives no information. The hypothesis of simple roots in
\cref{thm:split-simple} cannot be dropped.
\end{example}

\section{The spectral theorem in practice}

\begin{example}{Orthogonal diagonalisation of a symmetric matrix}{symmetric-example}
Let $A=\begin{pmatrix}2&1\\1&2\end{pmatrix}$, which is symmetric. Then
$\det(A-\lambda I)=(2-\lambda)^{2}-1=(\lambda-1)(\lambda-3)$, so
$\spec(A)=\set{1,3}$ with eigenvectors $(1,-1)$ and $(1,1)$. These are
orthogonal, as \cref{thm:spectral} predicts. Normalising,
\[
  Q=\frac{1}{\sqrt2}\begin{pmatrix}1&1\\-1&1\end{pmatrix},\qquad
  \transpose{Q}AQ=\diag(1,3).
\]
Since both eigenvalues are positive, $A$ is positive definite by
\cref{cor:spd-eigen}, and
$\tr A\cdot\tr A^{-1}=4\cdot\bigl(1+\tfrac13\bigr)=\tfrac{16}{3}\geq4=n^{2}$,
illustrating \cref{prob:trace-inequality}.
\end{example}

\section{Visual explanation: what rank measures}

\begin{center}
\begin{tikzpicture}[scale=0.95,>=Stealth,font=\small]
  % source space
  \draw[dmzrule] (0,0) rectangle (3,2.4);
  \node[text=dmzsoft,font=\footnotesize] at (1.5,2.65) {$E$, $\dim E=3$};
  \draw[dmzgreen,fill=dmzgreen!12] (0.5,0.4) rectangle (1.4,1.9);
  \node[text=dmzgreen,font=\footnotesize] at (0.95,0.15) {$\Ker u$};
  % target space
  \draw[dmzrule] (5.5,0) rectangle (8.5,2.4);
  \node[text=dmzsoft,font=\footnotesize] at (7,2.65) {$F$};
  \draw[dmztheorem,fill=dmztheorem!12] (6.1,0.9) rectangle (7.9,1.6);
  \node[text=dmztheorem,font=\footnotesize] at (7,0.6) {$\im u$, $\rank u=2$};
  % arrow
  \draw[->,dmzaccent,line width=0.9pt] (3.2,1.2) -- (5.3,1.2)
        node[midway,above,text=dmzaccent,font=\footnotesize] {$u$};
  \draw[->,dmzsoft,dashed] (1.0,1.15) -- (6.9,1.25);
\end{tikzpicture}
\end{center}

The kernel is the information destroyed by $u$; the image is the information
kept. \Cref{thm:rank-nullity} says that the two amounts add up to $\dim E$:
nothing is created and nothing disappears.

\chapter{Common Mistakes, Summary and Practice}
\label{ch:summary}

\section{Common mistakes}

\begin{mistakes}
\begin{itemize}
  \item \textbf{Writing $\Ker u=\emptyset$.} A kernel always contains $0$; the
        correct statement for an injective map is $\Ker u=\set{0}$.
  \item \textbf{Believing that $F\cup G$ is a subspace.} It is one only when one
        of the two contains the other (\cref{exo:union}).
  \item \textbf{Confusing free and generating.} A basis needs both; forgetting one
        half is the most frequent loss of marks.
  \item \textbf{Using rank--nullity in infinite dimension.}
        \Cref{thm:rank-nullity} assumes $\dim E<\infty$.
  \item \textbf{Deducing diagonalisability from a split annihilating polynomial
        with multiple roots.} The roots must be simple; see
        \cref{ex:jordan-block}.
  \item \textbf{Assuming $\rank(AB)=\rank A\cdot\rank B$.} Only the inequality
        of \cref{prop:rank-inequalities} holds.
  \item \textbf{Forgetting that eigenvalues depend on the field.} The matrix
        $\begin{pmatrix}0&-1\\1&0\end{pmatrix}$ has no eigenvalue over $\R$ and
        two over $\C$.
  \item \textbf{Adding dimensions of subspaces.} The correct formula is
        $\dim(F+G)=\dim F+\dim G-\dim(F\cap G)$.
\end{itemize}
\end{mistakes}

\section{Summary}

\begin{summarybox}
\begin{itemize}
  \item A vector space is a set where addition and scaling behave as expected;
        subspaces, sums and direct sums are the ways of building new ones.
  \item A basis is a system of unique addresses; its cardinality, the dimension,
        does not depend on the basis (\cref{thm:dimension-well-defined}).
  \item A linear map is determined by its values on a basis. Its kernel measures
        destroyed information, its image kept information, and
        $\dim E=\rank u+\dim\Ker u$ (\cref{thm:rank-nullity}).
  \item In equal finite dimensions, injective, surjective and bijective coincide
        (\cref{cor:iso-criterion}).
  \item Eigenspaces of distinct eigenvalues are independent
        (\cref{thm:eigen-independent}); $u$ is diagonalisable exactly when their
        dimensions add up to $\dim E$ (\cref{thm:diagonalisable-criterion}).
  \item A split annihilating polynomial with simple roots gives diagonalisability
        for free (\cref{thm:split-simple}).
  \item Every real symmetric matrix is orthogonally diagonalisable
        (\cref{thm:spectral}).
\end{itemize}
\end{summarybox}

\section{Practice exercises}

Solve these before reading \cref{ch:solutions}. Complete solutions are given
there.

\begin{exercise}{Polynomials of bounded degree}{poly-dim}
Show that $\K_{n}[X]=\set{P\in\K[X]\setst\deg P\leq n}$ is a vector space of
dimension $n+1$, and that $\bigl(1,X,\dots,X^{n}\bigr)$ is a basis.
\end{exercise}

\begin{exercise}{Union of two subspaces}{union}
Let $F,G$ be subspaces of $E$. Prove that $F\cup G$ is a subspace if and only if
$F\subseteq G$ or $G\subseteq F$.
\end{exercise}

\begin{exercise}{Rank of a sum}{rank-sum}
Let $u,v\in\Hom(E,F)$ with $\dim E<\infty$. Prove
$\abs{\rank u-\rank v}\leq\rank(u+v)\leq\rank u+\rank v$.
\end{exercise}

\begin{exercise}{Kernel of the square}{kernel-square}
Let $u\in\End(E)$ with $\dim E<\infty$. Show that
$\Ker u=\Ker u^{2}$ if and only if $\Ker u\cap\im u=\set{0}$, and that this is
also equivalent to $E=\Ker u\oplus\im u$.
\end{exercise}

\begin{exercise}{Projections}{projection-exo}
Let $p\in\End(E)$ with $p^{2}=p$ and $\dim E<\infty$. Show that
$E=\Ker p\oplus\im p$, that $\im p=\Ker(p-\id)$, and that $\tr p=\rank p$.
\end{exercise}

\begin{exercise}{Similar matrices}{similar}
Let $A,B\in\Mat{n}{\K}$ be similar, that is $B=P^{-1}AP$ for some invertible
$P$. Show that $A$ and $B$ have the same characteristic polynomial, the same
trace, the same determinant and the same rank. Show by an example that the
converse fails.
\end{exercise}

\section{Further reading}

\Textcite{axler2015} treats the reduction theory without determinants;
\textcite{gourdon2009} contains a large number of examination-style exercises of
the difficulty of \cref{ch:problems}.

\chapter{Doctoral Entrance Problems}
\label{ch:problems}

The problems below are of the type and difficulty of the linear algebra parts of
Algerian doctoral entrance examinations: a short statement, no hint, and a
complete proof expected in twenty to thirty minutes. Full solutions are in
\cref{ch:solutions}.
\newpage

\begin{problem}{(USTHB) — Épreuve N°01}{}
Soit l'ensemble $V = \{M \in \mathcal{M}_2(\mathbb{R}) \;/\; {}^tM = M\}$.
\begin{enumerate}
\item Montrer que $V$ est un sous-espace vectoriel de $\mathcal{M}_2(\mathbb{R})$.

\item Déterminer une base de $V$ ainsi que la dimension de $V$.

\end{enumerate}

On pose
$$
A=\begin{pmatrix}
1 & 1\\
0 & 2
\end{pmatrix}
$$
et pour toute matrice $M\in V$, on note
$$
f(M)=AM({}^tA).
$$

\begin{enumerate}
\setcounter{enumi}{2}

\item Montrer que $f$ est un endomorphisme de $V$.

\item Donner la matrice $B$ associée à $f$ dans la base de $V$ obtenue précédemment.

\item Déterminer le spectre de $f$.

\item L'endomorphisme $f$ est-il diagonalisable sur $\mathbb{R}$ ? Justifier.

\end{enumerate}

\source{\emph{\small Université des Sciences et de la Technologie Houari Boumediène (USTHB)}}
\end{problem}

\newpage

\begin{problem}{(USTHB) — Épreuve N°02}{antisymmetric-r3}
On munit $E=\mathbb{R}^3$ de sa base canonique $\mathcal{B}=\{e_1,e_2,e_3\}$ et de son produit scalaire usuel~:
pour tous $x=(x_1,x_2,x_3)$ et $y=(y_1,y_2,y_3)$ dans $E$,
$\langle x,y\rangle = x_1 y_1 + x_2 y_2 + x_3 y_3$.

Soit $u$ un endomorphisme non nul de $E$ tel que
$\langle u(x),x\rangle=0$ pour tout $x\in E$.
\begin{enumerate}
  \item Montrer que, pour tout $(x,y)\in E^{2}$,
        $\langle u(x),y\rangle = -\langle x,u(y)\rangle$.
        On pourra utiliser $z=x+y$ et calculer $\langle u(z),z\rangle$.
  \item \begin{enumerate}
          \item Démontrer que $0$ est la seule valeur propre réelle possible de $u$.
          \item Soit $P_u(X)$ le polynôme caractéristique de $u$. Quel est son degré~?
                Pourquoi a-t-il au moins une racine réelle~? En déduire que
                $P_u(X)=X\,R(X)$ où $R(X)$ est de degré~2 sans racines réelles.
          \item Montrer que $\ker(u)$ est de dimension~1. Quelle est la dimension de $\mathrm{Im}(u)$~?
        \end{enumerate}
  \item \begin{enumerate}
          \item Démontrer que si $x\in\ker(u)$ et $y\in\mathrm{Im}(u)$ alors
                $\langle x,y\rangle=0$.
          \item En déduire que $E$ est somme directe de $\ker(u)$ et $\mathrm{Im}(u)$.
        \end{enumerate}
\end{enumerate}

\source{\emph{\small Université des Sciences et de la Technologie Houari Boumediène (USTHB), 2026, Épreuve~n\textsuperscript{o}~02.}}
\end{problem}

\chapter{Fully Worked Solutions}
\label{ch:solutions}

\section{Solutions of the practice exercises}

\begin{solution}{\Cref{exo:poly-dim}}
$\K_{n}[X]$ contains $0$ and is stable under linear combinations, since
$\deg(\lambda P+\mu Q)\leq\max(\deg P,\deg Q)\leq n$; it is therefore a subspace
of $\K[X]$. The family $(1,X,\dots,X^{n})$ generates it by definition of a
polynomial of degree at most $n$. It is free: a relation
$\sum_{k=0}^{n}\lambda_kX^{k}=0$ is an equality of polynomials, so all
coefficients vanish. Hence it is a basis and $\dim\K_{n}[X]=n+1$.
\end{solution}

\begin{solution}{\Cref{exo:union}}
If $F\subseteq G$ then $F\cup G=G$ is a subspace, and symmetrically.

Conversely, assume $F\cup G$ is a subspace and that neither inclusion holds.
Choose $x\in F\setminus G$ and $y\in G\setminus F$. Then $x+y\in F\cup G$.
If $x+y\in F$, then $y=(x+y)-x\in F$, a contradiction; if $x+y\in G$, then
$x=(x+y)-y\in G$, again a contradiction. Hence one of the inclusions holds.
\end{solution}

\begin{solution}{\Cref{exo:rank-sum}}
The upper bound is \cref{prop:rank-inequalities}. For the lower bound, apply it
to $u=(u+v)+(-v)$:
\[
  \rank u\leq\rank(u+v)+\rank(-v)=\rank(u+v)+\rank v,
\]
so $\rank u-\rank v\leq\rank(u+v)$. Exchanging $u$ and $v$ gives
$\rank v-\rank u\leq\rank(u+v)$, and the two inequalities together give the
stated bound with the absolute value.
\end{solution}

\begin{solution}{\Cref{exo:kernel-square}}
The inclusion $\Ker u\subseteq\Ker u^{2}$ always holds.

\emph{Suppose $\Ker u=\Ker u^{2}$.} Let $x\in\Ker u\cap\im u$, say $x=u(y)$ and
$u(x)=0$. Then $u^{2}(y)=0$, so $y\in\Ker u^{2}=\Ker u$, hence $x=u(y)=0$.

\emph{Suppose $\Ker u\cap\im u=\set{0}$.} Let $x\in\Ker u^{2}$. Then
$u(x)\in\Ker u\cap\im u=\set{0}$, so $x\in\Ker u$.

Finally, in finite dimension, $\dim\Ker u+\rank u=\dim E$ by
\cref{thm:rank-nullity}; hence $\Ker u\cap\im u=\set{0}$ is equivalent to
$E=\Ker u\oplus\im u$, the dimensions already matching.
\end{solution}

\begin{solution}{\Cref{exo:projection-exo}}
For every $x\in E$ write $x=(x-p(x))+p(x)$. The first term lies in $\Ker p$
because $p(x-p(x))=p(x)-p^{2}(x)=0$, and the second lies in $\im p$; hence
$E=\Ker p+\im p$. If $x\in\Ker p\cap\im p$, say $x=p(y)$ with $p(x)=0$, then
$x=p(y)=p^{2}(y)=p(x)=0$, so the sum is direct.

If $x\in\im p$, $x=p(y)$, then $p(x)=p^{2}(y)=p(y)=x$, so $\im p\subseteq\Ker(p-\id)$;
the reverse inclusion is clear since $x=p(x)$ implies $x\in\im p$.

Choosing a basis of $\im p$ followed by a basis of $\Ker p$, the matrix of $p$ is
$\diag(1,\dots,1,0,\dots,0)$ with $r=\rank p$ ones. Hence $\tr p=r=\rank p$; the
trace does not depend on the basis by \cref{exo:similar}.
\end{solution}

\begin{solution}{\Cref{exo:similar}}
With $B=P^{-1}AP$,
\[
  \det(B-\lambda I)=\det\bigl(P^{-1}(A-\lambda I)P\bigr)
  =\det(P)^{-1}\det(A-\lambda I)\det(P)=\det(A-\lambda I),
\]
so the characteristic polynomials coincide, and with them the determinant
(value at $0$ up to sign) and the trace (coefficient of $\lambda^{n-1}$ up to
sign). Equality of ranks holds because $P$ and $P^{-1}$ are isomorphisms:
$\rank(P^{-1}AP)=\rank A$.

The converse fails: $I_2$ and $\begin{pmatrix}1&1\\0&1\end{pmatrix}$ have the
same characteristic polynomial $(X-1)^{2}$, the same trace, determinant and
rank, but they are not similar, since $I_2$ is diagonal while the second matrix
is not diagonalisable (\cref{ex:jordan-block} applied to $N=B-I_2$).
\end{solution}

\section{Solutions of the doctoral problems}

\begin{solution}{\Cref{prob:idempotent}}
\textbf{(1)} If $\lambda\in\spec(u)$ with eigenvector $x\neq0$, then
$u^{2}(x)=\lambda^{2}x$ and $u(x)=\lambda x$ give $\lambda^{2}=\lambda$, so
$\lambda\in\set{0,1}$. If $u=0$ then $\spec(u)=\set{0}$; if $u=\id$ then
$\spec(u)=\set{1}$; otherwise $\Ker u\neq\set{0}$ and
$\Ker(u-\id)=\im u\neq\set{0}$, so $\spec(u)=\set{0,1}$.

\textbf{(2)} The polynomial $X^{2}-X=X(X-1)$ annihilates $u$ and is split with
simple roots, so $u$ is diagonalisable by \cref{thm:split-simple} and
$E=\Ker u\oplus\Ker(u-\id)$. In a basis adapted to this decomposition, with
$r=\dim\Ker(u-\id)$ first,
\[
  \mathrm{Mat}(u)=\begin{pmatrix} I_r & 0\\ 0 & 0\end{pmatrix}.
\]

\textbf{(3)} From that matrix, $\tr u=r$ and $\rank u=r$, since
$\im u=\Ker(u-\id)$ by \cref{exo:projection-exo}. Hence $\tr u=\rank u$.

\textbf{(4)} $(\id-u)^{2}=\id-2u+u^{2}=\id-u$, so $\id-u$ is a projection. Its
image is $\Ker u$ and its kernel is $\im u$; therefore
$\rank u+\rank(\id-u)=\dim\im u+\dim\Ker u=n$ by \cref{thm:rank-nullity}.
\end{solution}

\begin{solution}{\Cref{prob:nilpotent}}
\textbf{(1)} If $Ax=\lambda x$ with $x\neq0$, then $0=A^{p}x=\lambda^{p}x$, so
$\lambda^{p}=0$ and $\lambda=0$; conversely $A$ is not invertible (otherwise
$A^{p}$ would be), so $0$ is an eigenvalue and $\spec(A)=\set{0}$. Over $\C$ the
characteristic polynomial is split, $\chi_A=\prod_{i}(X-\lambda_i)$ with all
$\lambda_i=0$, hence $\chi_A=X^{n}$.

\textbf{(2)} Since $A^{p}=0$,
\[
  (I_n-A)\bigl(I_n+A+\dots+A^{p-1}\bigr)=I_n-A^{p}=I_n,
\]
so $I_n-A$ is invertible with $(I_n-A)^{-1}=\sum_{k=0}^{p-1}A^{k}$.

\textbf{(3)} The eigenvalues of $I_n+A$ are $1+\lambda_i=1$ counted with
multiplicity, and the determinant is their product, so $\det(I_n+A)=1$. Likewise
$A^{k}$ is nilpotent, so all its eigenvalues are $0$ and
$\tr(A^{k})=\sum_i\lambda_i(A^{k})=0$.

\textbf{(4)} If $A$ is diagonalisable then $A$ is similar to $\diag(0,\dots,0)=0$
by (1), hence $A=0$. Conversely $A=0$ is diagonal.
\end{solution}

\begin{solution}{\Cref{prob:cubic}}
\textbf{(1)} The polynomial $X^{3}-X=X(X-1)(X+1)$ annihilates $u$, is split over
$\R$ and has simple roots, so \cref{thm:split-simple} gives
\[
  E=\Ker u\oplus\Ker(u-\id)\oplus\Ker(u+\id).
\]

\textbf{(2)} Concatenating bases of the three kernels gives a basis of
eigenvectors, so $u$ is diagonalisable and $\spec(u)$ is any nonempty subset of
$\set{-1,0,1}$; each subset is realised, for instance by a diagonal matrix with
the corresponding entries.

\textbf{(3)} In such a basis, $u$ has matrix $\diag(0,\dots,0,1,\dots,1,-1,\dots,-1)$.
Its rank is the number of nonzero diagonal entries, that is
$\dim\Ker(u-\id)+\dim\Ker(u+\id)$.

\textbf{(4)} For $\dim E=2$, the possible matrices up to similarity are the
diagonal matrices with entries in $\set{-1,0,1}$, namely
\[
  \diag(0,0),\quad\diag(1,0),\quad\diag(-1,0),\quad
  \diag(1,1),\quad\diag(-1,-1),\quad\diag(1,-1),
\]
six classes in all (the order of the diagonal entries does not change the
similarity class).
\end{solution}

\begin{solution}{\Cref{prob:trace-inequality}}
\textbf{(1)} By \cref{thm:spectral}, $A=QD\transpose{Q}$ with $Q$ orthogonal and
$D=\diag(\lambda_1,\dots,\lambda_n)$; by \cref{cor:spd-eigen} all
$\lambda_i>0$, so $\det A=\prod\lambda_i\neq0$ and
$A^{-1}=QD^{-1}\transpose{Q}$, which is symmetric with eigenvalues
$1/\lambda_i>0$, hence positive definite.

\textbf{(2)} Traces are invariant under similarity (\cref{exo:similar}), so
$\tr A=\sum_i\lambda_i$ and $\tr A^{-1}=\sum_i 1/\lambda_i$. By the
Cauchy--Schwarz inequality applied to the vectors
$\bigl(\sqrt{\lambda_i}\bigr)_i$ and $\bigl(1/\sqrt{\lambda_i}\bigr)_i$,
\[
  n^{2}=\Bigl(\sum_{i=1}^{n}\sqrt{\lambda_i}\cdot\frac{1}{\sqrt{\lambda_i}}\Bigr)^{2}
  \leq\Bigl(\sum_{i=1}^{n}\lambda_i\Bigr)\Bigl(\sum_{i=1}^{n}\frac{1}{\lambda_i}\Bigr)
  =\tr(A)\,\tr(A^{-1}).
\]

\textbf{(3)} Equality in Cauchy--Schwarz holds if and only if the two vectors are
proportional, that is $\lambda_i=c$ for all $i$ and some $c>0$. Then
$A=Q(cI_n)\transpose{Q}=cI_n$: equality holds exactly for the positive scalar
multiples of the identity.
\end{solution}

\begin{solution}{\Cref{prob:commuting}}
\textbf{(1)} Let $\lambda\in\spec(u)$ and $x\in E_\lambda$. Then
\[
  u\bigl(v(x)\bigr)=v\bigl(u(x)\bigr)=v(\lambda x)=\lambda v(x),
\]
so $v(x)\in E_\lambda$; thus $v(E_\lambda)\subseteq E_\lambda$.

\textbf{(2)} Since $u$ is diagonalisable, $E=\bigoplus_{\lambda}E_\lambda$. Fix
$\lambda$ and let $v_\lambda$ be the restriction of $v$ to $E_\lambda$, which is
well defined by (1). As $v$ is diagonalisable, it admits a split annihilating
polynomial with simple roots (its minimal polynomial); that polynomial also
annihilates $v_\lambda$, so $v_\lambda$ is diagonalisable by
\cref{thm:split-simple}. Choose a basis of $E_\lambda$ of eigenvectors of
$v_\lambda$; its vectors are also eigenvectors of $u$, with eigenvalue
$\lambda$. Concatenating over $\lambda\in\spec(u)$ gives a basis of $E$ in which
both $u$ and $v$ are diagonal.

\textbf{(3)} In that common basis, $u+v$ and $u\circ v$ are diagonal, hence
diagonalisable.
\end{solution}


%% ---------------- back matter ----------------
\backmatter
\printbibliography[heading=bibintoc,title={References}]
\printindex

\end{document}
